\section{Background}
\label{sec:background}

This section situates Recursive Spawning within the broader landscape of prior work across five domains: agent lifecycle management in multi-agent systems, accountability and provenance in distributed AI, fixed versus mutable delegation chains, hierarchical task decomposition in AI agents, and the BX3 Framework as the architectural substrate that makes immutable parent chains tractable to enforce at runtime.

% ------------------------------------------------------------------
\subsection{Agent Lifecycle Management in Multi-Agent Systems}

Multi-agent systems (MAS) composed of LLM-powered agents introduce a class of lifecycle challenges fundamentally different from those in conventional distributed systems. Each agent is itself a stateful, reasoning entity capable of planning, tool use, adaptive behavior, and recursive self-modification, making lifecycle management substantially more complex than traditional microservice orchestration \cite{alenezi2026}. Agent lifecycle in this context encompasses creation, authorization, execution, communication, suspension, and decommissioning of entities that may multiply in real time as orchestrator agents spawn sub-agents in response to emerging subtask demands \cite{moore2025hmas}.

The defining failure mode in unmanaged agent fleets is \textit{agent sprawl}: the proliferation of agents without clear ownership, explicit lifecycle boundaries, or consistent policy enforcement at handoff boundaries \cite{moore2025hmas}. This is not merely an operational inconvenience; it is an accountability liability. When an agent acts, the system must be capable of reconstructing \textit{which agent performed which action, under whose authorization, at what time, with what effect}. Without explicit lifecycle infrastructure, audit trails degenerate into siloed logs that cannot sustain regulatory scrutiny or forensic reconstruction \cite{provcagent2025, dzone2025accountability}.

Recommended controls include an identity and persona registry, orchestration and cross-domain mediation, runtime policy enforcement with kill-switch capability, and explicit decommissioning procedures that tie credential revocation to audit logging \cite{moore2025hmas, agentorchestra2025}. AgentOrchestra identifies a complementary concern: agents spawned for specialized subtasks frequently fail to inherit the policy context of their parent, creating enforcement gaps precisely at handoff boundaries \cite{agentorchestra2025}. PROV-AGENT extends the W3C PROV standard into agentic workflows, capturing agent-centric metadata (prompts, responses, decisions) and linking them to the broader workflow context, demonstrating that end-to-end provenance enables post-hoc fault isolation in heterogeneous multi-agent pipelines \cite{provcagent2025}. These works collectively establish that lifecycle management is not separable from accountability: you cannot have one without the other.

% ------------------------------------------------------------------
\subsection{Accountability and Provenance in Distributed Systems}

The accountability deficit in distributed AI systems differs in kind from conventional software accountability. When an LLM-powered agent makes a decision inside a multi-agent chain, the decision may emerge from complex, opaque reasoning over contextual state that is difficult to reconstruct post-hoc \cite{dzone2025accountability}. Without explicit provenance infrastructure, it is functionally impossible to attribute outcomes to responsible actors, verify policy compliance at each delegation hop, or detect scope drift when a sub-agent acts beyond its authorized mandate \cite{sentinelagent2026}.

SentinelAgent introduces the Delegation Chain Calculus with seven formal properties---authority narrowing, policy preservation, forensic reconstructibility, cascade containment, scope-action conformance, output schema conformance, and intent preservation---and demonstrates that enforcing these properties at runtime via a non-LLM Delegation Authority Service achieves 100\% true-positive rate on DelegationBench v4 across 516 scenarios \cite{sentinelagent2026}. Critically, properties P3--P7 are mechanically verified using TLA+ across 2.7 million states with zero violations \cite{sentinelagent2026}. The key insight for Recursive Spawning is that \textit{forensic reconstructibility} demands every delegation hop to be recorded immutably, not merely accessible via a mutable log.

The Human Delegation Provenance (HDP) Protocol provides a cryptographic implementation of this principle: a token binds a human principal to their declared authorization scope and session identifier, then each agent in the chain appends a cryptographically signed hop to an append-only chain, with offline verification requiring only the issuer's Ed25519 public key and the current session identifier---no central registry or third-party trust anchors \cite{hdp2025arxiv}. Chain of Consciousness extends this to an append-only SHA-256 hash chain of lifecycle events anchored to Bitcoin via OpenTimestamps, with W3C DID-based persistent agent identity enabling continuity proofs that link discontinuous sessions into a single verifiable identity arc \cite{chainofconsciousness2026}. TrustTrack and AIAuditTrack propose decentralized interaction graphs recorded on tamper-evident distributed ledgers, making agent-to-agent delegations auditable without relying on centralized log integrity \cite{trusttrack2025, aaudittrack2025}. Layered Mutability provides the conceptual model for distinguishing mutable from immutable elements in persistent agent systems, showing that selectively fixing certain structural elements while permitting controlled mutation elsewhere achieves a favorable tradeoff between adaptability and auditability \cite{layeredmutability2026}. Together, these works establish a clear consensus: provenance infrastructure is a prerequisite for accountable multi-agent deployment.

% ------------------------------------------------------------------
\subsection{Fixed versus Mutable Delegation Chains}

The tension between fixed and mutable delegation chains is a central unresolved question in agentic system design. A \textit{fixed delegation chain} establishes delegation relationships at spawn time and preserves them immutably throughout execution; the parent pointer cannot be severed or redirected without detection, and any attempt to do so constitutes a verifiable policy violation. A \textit{mutable delegation chain} permits dynamic re-delegation, agent migration, or role reassignment during execution, offering flexibility at the cost of forensic traceability.

Mutable delegation is the dominant approach in current LLM multi-agent frameworks. AgentSpawn adaptively recruits collaborators based on task decomposition but spawned agents carry no cryptographic record of their parent's identity beyond a shared conversation context \cite{agentspawn2025}. FIRMHIVE spawns hierarchical agent trees for firmware security analysis but relies on in-memory parent references lost on process restart \cite{firmhive2025}. In both cases, post-hoc accountability requires trusting the system's internal log integrity or reconstructing behavior from partial, mutable records. AgentOrchestra identifies that spawned agents may not inherit policy context from their parent, creating enforcement gaps at every handoff \cite{agentorchestra2025}. These are not implementation oversights; they are structural consequences of mutable delegation by design.

The HDP Protocol and SentinelAgent both independently conclude that mutable delegation is an attack surface exploited by prompt injection attacks that hijack in-flight delegations \cite{hdp2025arxiv, sentinelagent2026}. The consensus across these works is that immutable delegation with cryptographic identity binding is the required alternative for accountable systems. BX3's Immutable Parent Pointer resolves the fixed-versus-mutable tension architecturally: the pointer itself is immutable once established, enforced by the Fact Layer, while the content of what is delegated may be bounded, reinterpreted, and escalated through the delegation chain under controlled conditions defined by the Bounds Engine.

% ------------------------------------------------------------------
\subsection{Hierarchical Task Decomposition in AI}

Hierarchical task decomposition (HTD) addresses the fundamental scalability limit of flat agent architectures: a single agent cannot efficiently represent or execute the full complexity of a high-level objective over extended time horizons. HTD recursively breaks high-level objectives into subgoals handled by specialized agents, with a coordination layer managing execution order, data flow, and termination conditions \cite{reactree2025, deepagent2025, structuredagent2025}.

ReAcTree introduces hierarchical agent trees with two node types: control flow nodes that organize execution order, and agent nodes that reason, act, and expand into further subgoals \cite{reactree2025}. Dynamic tree expansion allows agent nodes to autonomously generate and refine subgoals at runtime, substantially outperforming flat sequential baselines (63\% versus 24\% goal success on WAH-NL using Qwen2.5 72B). StructuredAgent combines online hierarchical planning using dynamic AND/OR trees with a structured memory module that tracks candidate solutions and constraints, enabling efficient exploration of contingent action sequences when a path fails \cite{structuredagent2025}. DeepAgent uses a Hierarchical Task DAG (HTDAG) to dynamically decompose high-level objectives into interdependent sub-tasks across multiple layers, with a recursive two-stage planner-executor enabling continuous task refinement and real-time graph modification including user-directed intervention \cite{deepagent2025}. HiPlan introduces a two-level hierarchy where global milestone action guides provide high-level structure and subgoals while local step-wise hints adapt past trajectory segments to current observations, decoupling long-horizon structure from low-level adaptability \cite{hiplan2025}. ReCAP addresses context drift in purely sequential prompting and cross-level continuity loss in hierarchical prompting through plan-ahead decomposition, generating a full subtask list before execution begins \cite{recap2025}.

The Recursive Delegation Engine (RDE) provides a formal decision-theoretic treatment of HTD: agents use recursive decision rules to choose between executing, decomposing, or delegating subtasks, with success probabilities propagated back up the chain to inform parent decisions \cite{rde2025}. RDE's key insight is that the delegation graph itself is a learned structure---the optimal decomposition strategy adapts based on observed outcomes. However, RDE does not address the accountability implications of adaptive delegation decisions, treating them as optimization problems rather than as authorization events requiring provenance. The structural gap that Recursive Spawning fills is converting the HTD task graph into an accountable lineage graph by design: every spawned sub-agent is bound to its accountable ancestor with immutability guarantees that no HTD system currently provides as an architectural property.

% ------------------------------------------------------------------
\subsection{The BX3 Framework as Substrate for Immutable Parent Chains}

The BX3 Framework, as introduced in prior work, organizes any complex system into three immutable functional layers: the \textit{Purpose Layer} (intent, judgment, and upstream accountability), the \textit{Bounds Engine} (bounded reasoning, constrained execution, and controlled reinterpretation), and the \textit{Fact Layer} (deterministic enforcement, hard constraint, and forensic auditability) \cite{bx3framework}. The framework is actor-agnostic: any entity capable of satisfying a layer's functional requirements may occupy that layer---human, AI system, mechanical process, or hybrid composition. What cannot vary are the properties each layer must maintain.

The upstream accountability guarantee is the property most directly relevant to Recursive Spawning: when any node encounters a state it cannot resolve within its bounds, accountability escalates recursively upward through the system hierarchy, bypassing all machine actors, until it reaches a human anchor \cite{bx3framework}. The property \textit{fails upward into human consciousness---never downward into algorithmic chaos} is not a policy directive but an architectural constraint enforced by the Fact Layer. This guarantee applies at the system level; Recursive Spawning extends it to the \textit{per-node} level.

The BX3 Protocol's Immutable Parent Pointer converts the system-level accountability guarantee into a per-node architectural property. Every spawn event is treated as a delegation event with three formally defined components: a parent identity (immutable, enforced by the Fact Layer), a bounded mandate derived from the parent's Purpose Layer authorization (enforced by the Bounds Engine), and an immutable record in the Forensic Ledger (enforced by the Fact Layer). This converts the BX3 three-layer model from a system-level accountability architecture into a \textit{per-node accountability architecture}, enabling recursive spawning to preserve upstream accountability at every tier of the task decomposition hierarchy.

Critically, the Bounds Engine's bounded reinterpretation privilege ensures that a child node's Bounds Engine may reason within the mandate established by its parent, but it cannot expand that mandate unilaterally. If the child node encounters a state outside its bounded mandate, it escalates to its parent via the same upstream accountability mechanism. This creates a recursive accountability chain in which every node at every depth has a formally accountable parent, terminating at the human anchor at the root. The provenance of every decision is traceable to a human authorization event.

Recent complementary work confirms the timeliness of this contribution. SentinelAgent's Delegation Chain Calculus independently identifies forensic reconstructibility and scope-action conformance as required properties of accountable delegation \cite{sentinelagent2026}. PROV-AGENT demonstrates that W3C PROV semantics can be extended to capture agent-centric metadata in multi-agent workflows \cite{provcagent2025}. HDP provides the cryptographic primitives for append-only hop records \cite{hdp2025arxiv}. Layered Mutability provides the conceptual model for distinguishing mutable from immutable elements \cite{layeredmutability2026}. What none of these works supplies---and what the BX3 Framework with Recursive Spawning provides---is a unified, actor-agnostic architectural substrate that simultaneously enforces immutability of parent pointers, bounded reinterpretation by spawned nodes, and recursive escalation of unresolved states to a human anchor, all as first-class, formally specified properties rather than as implementation heuristics.

% ------------------------------------------------------------------
% REFERENCES
% ------------------------------------------------------------------
% NOTE TO AUTHOR: The following citations are drawn from the canonical
% bx3framework.bib file in the arxiv-upload directory. Ensure your
% .tex file uses \bibliography{bx3framework} (without .bib extension).
% All references below should be verified against the actual .bib file
% and removed or corrected if absent.

\end{document}